\chapter*{Resumo} 
\addcontentsline{toc}{chapter}{Resumo}

O desenvolvimento de sistemas de radar de onda contínua modulada em frequência (FMCW) requer como componente essencial um gerador de chirp, capaz de produzir um sinal cuja frequência varia linearmente no tempo. A qualidade deste gerador impacta diretamente na resolução de distância e velocidade do radar, condicionando significativamente seu desempenho para a medição de distâncias.

Neste trabalho apresenta-se o projeto, a implementação e a caracterização de um gerador de chirp para radares FMCW, destinado a gerar sinais de varredura em frequência com os requisitos solicitados para sua integração no sistema desenvolvido no Laboratório de Radiofrequência e Microondas (LARFyM). O projeto tem como objetivo sintetizar um sinal com uma largura de banda de varredura de 300 MHz — compreendida entre 1,55 GHz e 1,85 GHz — e um tempo de varredura de 1 ms para perfis em dente de serra e 0,5 ms para perfis triangulares. Para alcançar isso, emprega-se uma arquitetura baseada em um malha de captura de fase (PLL) utilizando o sintetizador de frequência ADF4159, um filtro de malha ativo com amplificador operacional, um oscilador controlado por tensão (VCO) e um divisor de potência Wilkinson, complementados por um filtro passa-faixa Hairpin no estágio de saída para suprimir componentes espúrias e harmônicas indesejadas.

Em uma primeira etapa, são realizados os cálculos teóricos e a modelagem de cada um dos blocos constitutivos, empregando ferramentas de simulação especializadas como ADIsimPLL e Advanced Design System (ADS). Posteriormente, realiza-se o projeto de placas de circuito impresso (PCB) com controle de impedância em tecnologia coplanar e microstrip, seguido pela fabricação, montagem e caracterização experimental de cada módulo separadamente. Finalmente, o sistema completo é integrado e seu funcionamento é validado através de medições de resposta em frequência, tensão de controle, linearidade da rampa e ruído de fase. Os resultados obtidos confirmam o correto desempenho do protótipo, validando sua integração no sistema de radar do LARFyM.

\vspace{1em}
\noindent\textbf{Áreas Temáticas do Projeto Integrador:} radiofrequência, micro-ondas, sistemas de comunicações e radar.

\noindent\textbf{Matérias:} Eletrônica Analógica III, Teoria do Campo Eletromagnético e Teoria das Comunicações. 

\noindent\textbf{Palavras-chave:} Radar FMCW, gerador de chirp, malha de captura de fase (PLL), oscilador controlado por tensão (VCO), filtro Hairpin.